\chapter{研究现状与相关工作}
\label{ch:survey}

本章沿“经典统计检测 $\to$ 大规模 MIMO 实用接收 $\to$ 迭代/消息传递 $\to$ 深度学习检测 $\to$ 核方法与稳健 CSI $\to$ 研究空白”的脉络梳理现状，并明确本工作的定位。

\section{经典多用户检测理论}

多用户检测的系统框架由 Verdú 等奠定\cite{verdú1998}：在高斯多址模型下，最优多用户检测、匹配滤波、解相关器与线性 MMSE 接收机构成清晰的性能--复杂度谱系。其核心洞见是：匹配滤波在多用户场景一般\textbf{不是}单用户最优的简单推广；利用干扰结构可显著降低误码。

对空间复用 MIMO，联合最大似然（ML）检测等价于有限星座上的最小距离搜索。球形译码（SD）及其软输入软输出变体通过树搜索逼近 ML，并可为信道译码器提供 LLR\cite{hassibi2005,hochwald2003,studer2011}。K-best、固定复杂度球译码等变体以可预测时延换取次优。然而当用户数 $K$ 与调制阶数同时增大（如 $K=40$、16-QAM），联合搜索即便经过剪枝，最坏情况复杂度仍难以满足基站实时处理约束；因而工业界上行接收长期以线性或轻度非线性算法为主。

\section{大规模 MIMO 下的线性与正则化接收}

大规模 MIMO 改变了检测问题的“典型工作点”\cite{marzetta2010,larsson2014,bjornson2017}：当 $M\gg K$ 时，信道硬化与渐近正交使匹配滤波/MF 与 MMSE 的性能迅速靠近最优\cite{hoydis2013}。实用算法谱系包括：
\begin{itemize}
  \item \textbf{MF / ZF / MMSE}：闭式线性估计加硬判决；MMSE 在干扰抑制与噪声放大之间折中，并可视为对角加载\cite{tse2005,kay1993}；
  \item \textbf{正则化 ZF 与加载因子选择}：用渐近确定性等价或仿真标定正则参数；
  \item \textbf{迭代求逆近似}：Neumann 级数、Gauss--Seidel、共轭梯度等，降低大规模矩阵求逆成本。
\end{itemize}
\textbf{MMSE+导频估计}因此成为可部署基线的“强标准答案”。本报告以其为主要对照，而非设立一个容易击败的弱基线。

需要强调的现状是：文献中“线性已足够”的结论高度依赖 $M/K$ 很大、SNR 不太低、CSI 质量较高等条件。当负荷比升高（本报告 $K/M=40/128$）或工作在 0--10\,dB 低中 SNR 时，线性接收与 MAP/ML 之间仍常存在可观 BER 间隙——这也是学习检测与非线性检测持续被研究的原因。

\section{干扰消除、格方法与近似消息传递}

在线性接收之上，研究界发展了多类增强：
\begin{enumerate}[leftmargin=2em]
  \item \textbf{SIC/PIC}：用硬/软符号重构干扰并消除；排序、可靠度加权影响误差传播；
  \item \textbf{格基约减 + 均衡}：改善等效信道条件数后再做 ZF/MMSE；
  \item \textbf{EP / AMP / GAMP / OAMP}：在大系统极限下具有状态演化分析，能逼近贝叶斯最优的某类可计算近似，并对相关信道有扩展（如 OAMP）\cite{he2018}。
\end{enumerate}
这些方法在\textbf{完美或高质量 CSI}下表现突出。它们与本工作的联系在于：高斯干扰近似、协方差白化等思想会出现在充分统计构造中；差异在于本项目最终决策落在 RKHS 假设类上最小化后验损失，而不是把 PIC/AMP 迭代本身作为主检测器。项目早期曾试验 residual/PIC 前端，实验上未成为优于 $z_{\mathrm{rob}}$+RKHS 的默认可部署路径。

\section{深度学习检测：数据驱动与模型展开}

物理层深度学习综述表明，上行大规模 MIMO 检测已成为 DL 应用的热点之一\cite{o2017}。按架构大致可分为两类\cite{samuel2019,he2018,khani2020}。

\subsection{数据驱动（黑盒）网络}
以 $\mathbf{y}$ 或 $(\mathbf{y},\hat{\mathbf{H}})$ 为输入，多层网络直接输出符号/比特软信息。优点是表达力强、可端到端训练；缺点是参数量大、可解释性弱、对信道分布偏移敏感，且多数工作默认较理想的 CSI。

\subsection{模型驱动展开（deep unfolding）}
将经典迭代检测的每一步展开为网络层，只学习少量参数。代表性工作包括：
\begin{itemize}
  \item \textbf{DetNet}：展开投影梯度型迭代，展示 DL 做 MIMO 检测的可行性\cite{samuel2019}；
  \item \textbf{OAMP-Net / OAMP-Net2}：展开正交近似消息传递，参数更少，并讨论与信道估计结合的不完美 CSI 情形\cite{he2018}；
  \item \textbf{MMNet 等}：面向相关信道的展开检测，强调经典线性方法在相关信道上的不足\cite{khani2020}。
\end{itemize}
现状共识可概括为：（i）展开往往比纯黑盒更稳、更省参；（ii）\textbf{底层迭代算法的选择常常比“多加几个可学习标量”更关键}；（iii）许多工作仍在相对较小的 $M,K$ 或较强 CSI 假设下评估；（iv）不完美 CSI、导频开销与跨 SNR 泛化仍是开放问题。

本报告用 CNN 作为深度学习对照，而不是复现全部 DetNet/OAMP-Net 族：目的是提供“数据驱动分类器”参照，主贡献放在\textbf{贝叶斯损失 + 结构化特征 + RKHS 决策形式}这一不同于展开网络的路线。

\section{核方法、多核学习及其在通信中的位置}

核方法通过再生核将线性算法隐式提升到高维特征空间\cite{schafer2006,shawe2004,vapnik1998}。高斯核等通用核在适当条件下具有万能逼近性质\cite{steinwart2008}；多核学习（MKL）在基核凸组合上学习权重，减轻核选择困难\cite{bach2004,micchelli2005,cortes2009,rakotomamonjy2008}。

在通信与信号处理中，核自适应滤波、高斯过程、核贝叶斯多用于信道预测、非线性均衡与调制识别等。直接把核分类器用在大规模 MU-MIMO 符号检测的系统工作相对少于深度学习。现有思路的常见不足是：
\begin{itemize}
  \item 输入仍用原始高维 $\mathbf{y}$，使后验在高 SNR 下极尖锐，核带宽难选；
  \item 损失与贝叶斯检测准则对齐不清晰；
  \item 未系统区分“特权信息上界”与“可部署信息集”。
\end{itemize}
本工作针对这些不足：以 16 维 GA/稳健软分为坐标、以 $J(f)$ 为训练目标、以 Adaptive-MKL/核内系数学习增强表达，并显式分开 Oracle 与可部署主线。

\section{不完美 CSI 与稳健接收现状}

不完美 CSI 下的稳健接收是经典课题\cite{kay1993,bjornson2017}：对角加载、最坏情况设计、贝叶斯估计--检测联合、以及把误差协方差并入干扰模型等。TDD 大规模 MIMO 还有导频污染问题\cite{marzetta2010}。深度学习检测文献也开始讨论联合信道估计与检测（JCESD）及估计误差感知训练\cite{he2018}，但总体仍少于完美 CSI 设定下的结果。

本报告采取可实现的中间路线：不宣称解决导频污染，而用导频 LS 误差方差 $\sigma_e^2$ 对干扰协方差做对角加载，得到 $z_{\mathrm{rob}}$。这与“加长导频硬追完美 CSI”不同——后者把增益押在估计质量上；我们把增益押在\textbf{检测器对估计误差的容忍与后验损失优化}上。

\section{评价准则与实验惯例}

文献中常见指标包括 SER/BER、与 ML 的差距、互信息/软输出质量、复杂度（FLOPs、时延）、以及跨信道相关模型/跨 SNR 的泛化。深度学习工作还常报告训练样本需求与在线推理时间。本报告主报：
\begin{itemize}
  \item 期望用户 bit BER（相对 MMSE 的增益）；
  \item 后验损失 $J$；
  \item Oracle 相对 MLD 的 BER/$J$/MSE（表达力上界）。
\end{itemize}
复杂度的严格实时基准留作后续工作。

\section{研究空白与本工作定位}

综合现状，可见如下空白与张力：
\begin{enumerate}[leftmargin=2em]
  \item 在 $K/M$ 不可忽略、低中 SNR、仅可部署 CSI 时，相对 MMSE 的\textbf{稳定可复现增益}仍有需求；
  \item 深度展开很强，但与“\textbf{保留表示定理型决策形式、直接最小化贝叶斯后验损失}”的核路线之间，系统对比与协同仍不足；
  \item 上界实验（真 $H$、$f^*$）与上线方案（$\hat H$、hard）常被混用，导致“可逼近性”与“可部署性”论证不清。
\end{enumerate}

\noindent\textbf{本工作定位：}在方案 A（只检 $x_1$）下，给出一条结构化 RKHS 检测主线——可部署时用 $z_{\mathrm{rob}}(\hat H)$+hard 追击 MMSE；Oracle 时用 $z(H)$+$f^*$ 验证核空间可逼近后验。它介于经典线性接收与黑盒/展开深度检测之间，强调\textbf{模型启发特征、贝叶斯准则与可分析假设类}的统一。
